\section{Kohn--Sham calculations in practice}\label{app:E.1}

A density-functional calculation of silicon at its measured lattice constant returns, in seconds on a few processor cores, a total energy, an electron density and a set of Kohn--Sham bands with an indirect gap near \SI{0.6}{eV}, against the \SI{1.17}{eV} measured at low temperature~\cite[Eq.~(11)]{ramanathan2020ionization}. The work of a calculation is knowing which outputs to trust, how far, and what each numerical parameter does to them. The calculation of silicon used here has the PBE functional, the norm-conserving silicon pseudopotential of the SG15 library~\cite{schlipf2015optimization} with four valence electrons per atom, and the diamond structure with two atoms per cell.

\subsection{Computed quantities}

The Kohn--Sham equations \cref{eq:3.4-KS} are solved for the periodic parts $u_{n\vect{k}}$ of the Bloch states on a grid of crystal momenta. The states define the density
\begin{equation}
  \nd(\vect{r}) = \frac{2}{n_k}\sum_{\vect{k}}\sum_{n}f_{n\vect{k}}\,|u_{n\vect{k}}(\vect{r})|^2 ,
  \label{eq:E.1-density}
\end{equation}
in which the factor 2 counts spin, $n_k$ is the number of grid points and $f_{n\vect{k}}$ are the occupations. Since $u_{n\vect{k}}$ is normalized over the cell (\cref{app:A.5}), $\nd$ integrates over $\Omega$ to the number of valence electrons per cell. The density fixes the Kohn--Sham potential $V_{\mathrm{KS}}$, the potential fixes the states, and the loop is iterated until the input and output densities agree. The total energy, the density, and the forces and stresses derived from the energy are ground-state quantities, exact in principle for the exact functional (\cref{sec:3.4}); their errors come from the approximate functional and from the numerics. The eigenvalues $\varepsilon_{n\vect{k}}$ and states $\psi_{n\vect{k}}$ are auxiliary. Apart from the highest occupied level, they have no exact physical meaning, and the band gap they give is too small even for the exact functional (\cref{eq:3.4-gap}). The response quantities built from the states, such as $\varepsilon_\infty$, the phonons or the density vertex $\rho_{mn}(\vect{k}, \vect{q})$, inherit the errors of the ground-state quantities and the errors of the eigenvalues and states.

\subsection{Functionals}

A calculation chooses one of the approximations to the exchange--correlation energy $E_{\mathrm{xc}}$ of \cref{sec:3.4}. PBE~\cite{perdew1996generalized} depends on the density gradient through $|\nabla\nd|/\nd^{4/3}$, and its form is built to satisfy exact constraints of the electron gas. In the screened hybrid HSE~\cite{heyd2003hybrid} the Coulomb kernel $e^2/4\pi r$ is split, with an error function, into a part that decays on a length of a few \AA{} and a long-range remainder, and the exact exchange acts with the short-range part only. Cutting off the long-range part models a semiconductor, where the exchange interaction between distant electrons is screened by the other electrons (\cref{sec:4.8}).

LDA binds too strongly and gives lattice constants that are slightly too small, by \SIrange{0.2}{0.6}{\percent} for the materials of \cref{tab:E-functionals}; PBE binds too weakly and gives them too large, by \SIrange{0.9}{2.1}{\percent}. Over the elemental crystals, PBE equilibrium volumes are too large by \SI{3.6}{\percent} on average~\cite[Table~IV]{lejaeghere2014error}, about \SI{1.2}{\percent} in the lattice constant. The semilocal functionals underestimate gaps, LDA and PBE alike, and close the small gap of germanium altogether; HSE recovers most of the measured gap. Against the measured gaps of 41 semiconductors and insulators, HSE has a mean absolute error of \SI{0.38}{eV}, falling to \SI{0.18}{eV} for the typical semiconductors of the set, whose gaps are measured at low temperature~\cite[Table~I and text]{garza2016predicting}. \Cref{tab:E-functionals} collects the numbers for silicon, germanium and GaAs. In silicon the hybrid lattice constant is \SI{0.6}{\percent} too large and the hybrid gap \SI{0.05}{eV} too large. In germanium, where LDA gives a metal, the hybrid reopens the gap but leaves it \SI{0.2}{eV} too small.

\begin{table}[t]
  \centering
  \caption{Lattice constants and band gaps of silicon, germanium and gallium arsenide for three functionals, without spin--orbit coupling. Lattice constants from LDA and PBE are all-electron values from Ref.~\cite[Table~II]{csonka2009assessing}, those from HSE from Ref.~\cite[Table~I]{lucero2012improved}. Gaps are evaluated at each functional's own relaxed lattice constant, from Ref.~\cite[Table~I]{lucero2012improved}, whose LDA lattice constant of silicon is \SI{5.410}{\angstrom}; a zero gap means that the functional makes the crystal a metal. The static-lattice constants are the measured ones extrapolated to zero temperature with the zero-point anharmonic expansion removed~\cite[Table~II]{csonka2009assessing}; they are not measured at any temperature, and they lie below the lattice constants measured at low temperature. They are used only inside calculations, as the cell of a static lattice. The room-temperature lattice constants are measured: silicon~\cite[lattice parameter of silicon]{mohr2025codata}, germanium~\cite[Sect.~II\,A]{rideau2006strained} and GaAs~\cite{vurgaftman2001band}; they are the values for observables and the cell of the shared lattice setting of \cref{app:E.17}. The measured gaps are low-temperature values~\cite[Table~I]{lucero2012improved}.}
  \label{tab:E-functionals}
  \begin{tabular}{l ccccc ccc}
    \toprule
    & \multicolumn{5}{c}{lattice constant (\AA)} & \multicolumn{3}{c}{gap (eV)} \\
    \cmidrule(lr){2-6}\cmidrule(lr){7-9}
    & LDA & PBE & HSE & static lattice & room $T$ (expt.) & LDA & HSE & measured \\
    \midrule
    Si   & 5.403 & 5.466 & 5.451 & 5.416 & 5.431 & 0.51 & 1.22 & 1.17 \\
    Ge   & 5.623 & 5.759 & 5.706 & 5.640 & 5.658 & 0.00 & 0.54 & 0.74 \\
    GaAs & 5.605 & 5.746 & 5.711 & 5.638 & 5.653 & 0.43 & 1.18 & 1.52 \\
    \bottomrule
  \end{tabular}
\end{table}

A hybrid is expensive because of the exact exchange. A semilocal potential acts on a state $\psi_{n\vect{k}}$ by a multiplication in real space. The Fock operator instead needs, for every pair of the state with an occupied state $\psi_{m\vect{k}'}$, the potential of the pair density $\psi^*_{m\vect{k}'}\psi_{n\vect{k}}$. Each pair potential costs one Poisson solve: a fast Fourier transform (FFT) of the pair density, a division by $|\vect{k} - \vect{k}' + \vect{G}|^2$ and an FFT back. With $N_{\mathrm{occ}}$ occupied bands, applying the Fock operator to all $N_{\mathrm{b}}$ bands at one $\vect{k}$ therefore costs about $n_kN_{\mathrm{occ}}N_{\mathrm{b}}$ FFTs, against $N_{\mathrm{b}}$ for a semilocal potential, a factor $n_kN_{\mathrm{occ}}$ more. For silicon on an $8^3$ grid this factor is $512\times4 \approx 2\times10^3$. In practice the inner sum over $\vect{k}'$ runs over a coarser grid than the outer one, which reduces the factor by one or two orders of magnitude. A hybrid calculation of a two-atom cell therefore costs $10$ to $10^3$ times a semilocal one, minutes to hours where PBE takes seconds.

In this book the choice of functional is usually PBE (or LDA) for structures, phonons and electron--phonon couplings, where semilocal functionals are accurate to a few percent. Wherever an excitation energy sets a threshold, a gap correction, a hybrid functional or a quasiparticle calculation is used instead. The two choices do not mix freely: a phonon computed with PBE and a threshold computed with HSE describe two slightly different crystals, and the lattice constant at which both are evaluated must be stated.

\subsection{Plane waves and the cutoff}

The periodic part of each Bloch state is expanded in plane waves,
\begin{equation*}
  u_{n\vect{k}}(\vect{r}) = \Omega^{-1/2}\sum_{\vect{G}}c_{n\vect{k}}(\vect{G})\,\ee^{\ii\vect{G}\cdot\vect{r}}, \qquad \frac{|\vect{k} + \vect{G}|^2}{2m_e} \le E_{\mathrm{cut}} ,
\end{equation*}
truncated at the kinetic-energy cutoff $E_{\mathrm{cut}}$, the one parameter of the basis. Codes quote $E_{\mathrm{cut}}$ in rydberg ($\SI{1}{Ry} = \SI{13.606}{eV}$) or hartree (two rydberg). The number of plane waves follows from counting reciprocal-lattice points in a sphere of radius $k_{\mathrm{cut}} = (2m_eE_{\mathrm{cut}})^{1/2}$. Each point occupies a volume $(2\pi)^3/\Omega$ of reciprocal space, so
\begin{equation}
  N_{\mathrm{PW}} \simeq \frac{4\pi k_{\mathrm{cut}}^3/3}{(2\pi)^3/\Omega} = \frac{\Omega\,k_{\mathrm{cut}}^3}{6\pi^2} .
  \label{eq:E.1-npw}
\end{equation}
For silicon, $\Omega = \SI{40.05}{\angstrom^3}$ (\cref{sec:3.1}) is 270.3 cubic bohr. At $E_{\mathrm{cut}} = \SI{60}{Ry}$, with $\hbar^2/2m_e = \SI{1}{Ry.bohr^2}$, $k_{\mathrm{cut}} = \sqrt{60}$ inverse bohr and \cref{eq:E.1-npw} gives 2121 plane waves per state. The density is a product of two states and has Fourier components up to $2k_{\mathrm{cut}}$, eight times as many as the plane waves of one state, about $1.7\times10^4$. These density components set the size of the FFT grid on which the potential is applied. A component $\vect{G}$ has integer coordinates $m_i = \vect{G}\cdot\vect{a}_i/2\pi$ with $|m_i| \le 2k_{\mathrm{cut}}|\vect{a}_i|/2\pi$; for $|\vect{a}_i| = a/\sqrt{2}$ the bound is 17.9. The indices then run from $-17$ to 17, so the FFT grid needs at least 35 points along each primitive vector. In practice the grid has 36, the next size that the FFT handles efficiently, and $36^3 \approx 4.7\times10^4$ points in all.

The cutoff is a variational parameter. For a fixed pseudopotential, enlarging the basis can only lower the total energy, so the energy converges from above and monotonically. The change in total energy between successive cutoffs estimates the remaining error only once a cutoff scan has reached its asymptotic regime: in germanium the change from 50 to \SI{60}{Ry}, \SI{1.3}{meV}, is less than half the error that remains at \SI{60}{Ry}, \SI{3.1}{meV}. \Cref{tab:E-cutoff} shows the scan for silicon and germanium. Silicon converges to \SI{1}{meV} per atom at \SI{50}{Ry}, and to \SI{0.1}{meV} at about \SI{60}{Ry}. Germanium, whose pseudopotential carries the ten 3d electrons in the valence, needs \SI{80}{Ry} for \SI{0.2}{meV}, and at \SI{30}{Ry} the germanium total energy is wrong by almost \SI{3}{eV} per atom. The pressure behaves differently in the two cases: in silicon it is stable to \SI{0.3}{kbar} from \SI{30}{Ry}, while in germanium it is not even of the right sign at \SI{30}{Ry}. The positive pressures at the measured lattice constants, \SI{24}{kbar} and \SI{38}{kbar}, show the PBE tendency to expand the lattice.

\begin{table}[t]
  \centering
  \caption{Convergence of the total energy and pressure with the plane-wave cutoff in the representative PBE calculations of silicon and germanium at their measured room-temperature lattice constants, \SI{5.431}{\angstrom} and \SI{5.658}{\angstrom} (\cref{tab:E-functionals}). The norm-conserving pseudopotentials carry four (Si) and fourteen (Ge, with 3d) valence electrons, and the grid has $8^3$ points $\vect{k}$. $\Delta E$ is the total energy per atom minus its value at \SI{120}{Ry}.}
  \label{tab:E-cutoff}
  \begin{tabular}{r rr rr}
    \toprule
    & \multicolumn{2}{c}{Si} & \multicolumn{2}{c}{Ge} \\
    \cmidrule(lr){2-3}\cmidrule(lr){4-5}
    $E_{\mathrm{cut}}$ (Ry) & $\Delta E$ (meV) & $P$ (kbar) & $\Delta E$ (meV) & $P$ (kbar) \\
    \midrule
    30 & 8.36 & 23.6 & 2880 & $-1615$ \\
    40 & 2.52 & 23.4 & 11.7 & 27.5 \\
    50 & 0.71 & 23.6 & 4.42 & 37.6 \\
    60 & 0.14 & 23.7 & 3.10 & 37.2 \\
    70 & 0.02 & 23.7 & 0.49 & 37.7 \\
    80 & 0.01 & 23.7 & 0.22 & 38.0 \\
    100 & 0.00 & 23.7 & 0.04 & 38.0 \\
    \bottomrule
  \end{tabular}
\end{table}

Different quantities converge with the cutoff at different rates, and the stress picks up a Pulay term. Energy differences between two structures converge faster than energies, since most of the cutoff error comes from the region near the nuclei and is the same in both structures. Absolute quantities, and quantities that probe the region near the nuclei, converge more slowly. The density vertex at large momentum transfer is the extreme case, with the cutoff condition \cref{eq:3.5-Ecut}. At a fixed cutoff the number of plane waves changes with the cell volume, so the basis is not the same at two lattice constants, and the computed stress carries the Pulay term, a spurious negative contribution that vanishes only at convergence. A structural relaxation must therefore be done at a cutoff where the stress, and not only the energy, is converged.

\subsection{Brillouin-zone sampling}

Sums over the zone become averages over a uniform grid of $n_k = n_1n_2n_3$ points, possibly shifted from the origin as proposed by Monkhorst and Pack~\cite{monkhorst1976special}. \Cref{app:C.2} derived the error of a uniform grid, \cref{eq:C.2-grid}: it is the sum of the real-space Fourier components of the integrand at the superlattice vectors $\vect{R}$ of the $n_1\times n_2\times n_3$ block of cells that the grid implies. For an insulator these Fourier components decay exponentially, so a modest grid suffices. In the representative silicon calculation at \SI{60}{Ry}, the total energy changes by \SI{1.6}{meV} per atom from $8^3$ to $10^3$, by \SI{0.3}{meV} from $10^3$ to $12^3$ and by \SI{0.06}{meV} from $12^3$ to $14^3$. Symmetry reduces the work: on a $\Gamma$-centered $12^3$ grid, the 48 operations of the diamond structure leave 72 inequivalent points.

The grid needed depends on the quantity, however, and a converged total energy says little about a converged response. The electronic dielectric constant $\varepsilon_\infty$ (\cref{sec:4.4}), computed with density-functional perturbation theory at the PBE lattice constant, is an example. It is 14.36 on an $8^3$ grid, 13.57 on $10^3$, 13.29 on $12^3$ and 13.20 on $14^3$; at the measured lattice constant it reaches 13.04 on $18^3$. The total energy has converged to \SI{1}{meV} per atom long before $\varepsilon_\infty$ converges to one percent. The reason is that a response involves transitions between occupied and empty states, each weighted by the factor $1/(\varepsilon_{c\vect{k}} - \varepsilon_{v\vect{k}})$. These factors are largest near the minimum direct gap, and a grid must resolve the region of the zone where they are largest. The same holds, more severely, for any quantity that probes a small region of the zone, such as the conduction-band minimum of silicon at $0.84$ of the way to X or the valley edges of a phonon-assisted window.

\subsection{Metals and smearing}

In a metal the occupations jump from one to zero at the Fermi surface. The integrand of \cref{eq:C.2-grid} then has a discontinuity, its Fourier components decay only as a power of $|\vect{R}|$, and a grid converges slowly and erratically. The remedy is to replace the step by a smooth function of width $\sigma_{\mathrm{sm}}$, $f_{n\vect{k}} = \tilde f\bigl((\mu_{\mathrm{ch}} - \varepsilon_{n\vect{k}})/\sigma_{\mathrm{sm}}\bigr)$, which makes the integrand smooth on the scale of the grid. The smearing changes the computed quantities, and the change must be controlled. The simplest case is Fermi--Dirac smearing, $\tilde f(x) = 1/(\ee^{-x} + 1)$, where $\sigma_{\mathrm{sm}}$ is a fictitious temperature. The quantity that a self-consistent calculation minimizes is then the free energy $F_{\mathrm{sm}} = E - \sigma_{\mathrm{sm}}S_{\mathrm{sm}}$, with the entropy $S_{\mathrm{sm}}$ of the occupations. Let the density of states be smooth on the scale of $\sigma_{\mathrm{sm}}$, with $\mathcal{N}_{\mathrm{F}}$ per spin and volume at the Fermi level (\cref{sec:2.5}), so $2\Omega\mathcal{N}_{\mathrm{F}}$ per cell. The Sommerfeld expansion then gives
\begin{equation}
  E(\sigma_{\mathrm{sm}}) = E(\sigma_{\mathrm{sm}} = 0) + \frac{\pi^2}{3}\,\Omega\mathcal{N}_{\mathrm{F}}\,\sigma_{\mathrm{sm}}^2, \qquad
  F_{\mathrm{sm}}(\sigma_{\mathrm{sm}}) = E(\sigma_{\mathrm{sm}} = 0) - \frac{\pi^2}{3}\,\Omega\mathcal{N}_{\mathrm{F}}\,\sigma_{\mathrm{sm}}^2 ,
  \label{eq:E.1-sommerfeld}
\end{equation}
up to terms of order $\sigma_{\mathrm{sm}}^4$, with $E(\sigma_{\mathrm{sm}} = 0)$ the energy per cell of the unsmeared metal. The first relation is the electronic heat of a metal, with $\sigma_{\mathrm{sm}}$ in the role of the temperature; the second follows from the electronic entropy at the same order, $S_{\mathrm{sm}} = (2\pi^2/3)\Omega\mathcal{N}_{\mathrm{F}}\sigma_{\mathrm{sm}}$. The errors of $E$ and $F_{\mathrm{sm}}$ are equal and opposite, and the average $\tfrac12(E + F_{\mathrm{sm}})$ is correct to order $\sigma_{\mathrm{sm}}^4$. The forces, however, are derivatives of $F_{\mathrm{sm}}$, not of $E$, so they carry the full $\sigma_{\mathrm{sm}}^2$ error.

Other smearing functions remove the $\sigma_{\mathrm{sm}}^2$ term from the start. The functions of Methfessel and Paxton~\cite{methfessel1989high} are Gaussians multiplied by Hermite polynomials. The low moments of these smearing functions vanish, but the occupations they give are not confined to $[0, 1]$. The ``cold smearing'' of Marzari and Vanderbilt~\cite[Eq.~(1)]{marzari1999thermal},
\begin{equation}
  \frac{\dd\tilde f}{\dd x} = \frac{1}{\sqrt{\pi}}\,\exp\!\Bigl[-\Bigl(x - \frac{1}{\sqrt{2}}\Bigr)^2\Bigr]\bigl(2 - \sqrt{2}\,x\bigr),
  \label{eq:E.1-cold}
\end{equation}
is the broadening that replaces the delta function. The occupations obtained by integrating this broadening are never negative, although they can exceed one. The broadening is chosen so that the free energy has no quadratic dependence on $\sigma_{\mathrm{sm}}$, and neither do the derivatives of the free energy, forces and stresses included~\cite{marzari1999thermal}. A width of a few tenths of an eV is common for a structural calculation of a simple metal. The smearing width $\sigma_{\mathrm{sm}}$ must not be confused with a physical temperature. For a quantity that depends on the occupations near the Fermi surface, such as the density of states at $\mu_{\mathrm{ch}}$ of aluminum (\cref{sec:7.2}) or a Fermi-surface average of the electron--phonon coupling, the smearing and the grid must be converged together. The tetrahedron method of \cref{app:C.2}~\cite{blochl1994improved} avoids smearing altogether. An insulator such as silicon needs no smearing: its occupations are fixed at one and zero, and smearing an insulator is an error.

\subsection{Self-consistency}

The density that goes into $V_{\mathrm{KS}}$ and the density that comes out of the states must agree. A calculation starts from a guess, usually a superposition of atomic densities, and iterates. Feeding the output density straight back in diverges in all but the smallest systems, and the reason is screening. Let the input density differ from the self-consistent density by the input error $\delta\nd_{\mathrm{in}}$. The Hartree potential of the input error, $v_{\mathrm{C}}\,\delta\nd_{\mathrm{in}}$, changes the states. The output density then differs from the self-consistent density by the output error $\delta\nd_{\mathrm{out}} = R^0v_{\mathrm{C}}\,\delta\nd_{\mathrm{in}}$, with $R^0$ the static independent-particle response of \cref{sec:4.6} (we drop the exchange--correlation kernel, which does not change the argument). The dielectric matrix of the random-phase approximation is $\varepsilon = 1 - v_{\mathrm{C}}R^0$ (\cref{sec:4.4}); the same product in the other order, $\varepsilon^{\mathrm{T}} = 1 - R^0v_{\mathrm{C}}$, acts on densities and has the same eigenvalues. The output error is $(1 - \varepsilon^{\mathrm{T}})\,\delta\nd_{\mathrm{in}}$. Linear mixing takes the new input as $\nd_{\mathrm{in}} + \beta_{\mathrm{mix}}(\nd_{\mathrm{out}} - \nd_{\mathrm{in}})$, and the input error then evolves as
\begin{equation}
  \delta\nd_{\mathrm{in}}^{(i+1)} = (1 - \beta_{\mathrm{mix}}\,\varepsilon^{\mathrm{T}})\,\delta\nd_{\mathrm{in}}^{(i)} .
  \label{eq:E.1-mixing}
\end{equation}
The iteration converges if $|1 - \beta_{\mathrm{mix}}\varepsilon_\lambda| < 1$ for every eigenvalue $\varepsilon_\lambda$ of the static dielectric matrix, that is for $\beta_{\mathrm{mix}} < 2/\max_\lambda\varepsilon_\lambda$. In an insulator the largest eigenvalue stays finite as the cell grows, of the order of $\varepsilon_\infty$, and a fixed mixing parameter works. In a metal, the Thomas--Fermi form $\varepsilon(\vect{G}) = 1 + k_{\mathrm{TF}}^2/G^2$ (\cref{sec:4.4}) grows without limit at small $G$, and the smallest $G$ of a cell of length $L$ is $2\pi/L$. A large metallic cell therefore needs $\beta_{\mathrm{mix}} \lesssim 2(2\pi/k_{\mathrm{TF}}L)^2$, and long-wavelength charge oscillates back and forth between iterations, a failure known as charge sloshing. The cure of Kerker~\cite{kerker1981efficient} is to mix each Fourier component with $\beta_{\mathrm{mix}}G^2/(G^2 + k_{\mathrm{TF}}^2)$, with a wavevector of the order of $k_{\mathrm{TF}}$ in the denominator. The Kerker factor cancels the $1/G^2$ of the Thomas--Fermi dielectric function, so that $\beta_{\mathrm{mix}}\varepsilon(\vect{G}) \approx \beta_{\mathrm{mix}}$ for all $\vect{G}$. Production codes combine a preconditioner such as Kerker's with a quasi-Newton update that uses the history of previous iterations to build an approximate inverse of $\varepsilon$, the modified Broyden method of Johnson~\cite{johnson1988modified} or the closely related Pulay mixing.

The convergence threshold of the self-consistency loop must match the use of the result. A total energy to \SI{0.1}{meV} needs a loose threshold, since the energy is variational and its error is quadratic in the density error. A response calculation started from the converged ground state needs a tight one, since the first-order change of the states is computed against the ground-state Hamiltonian, and a residual error in the density appears linearly in the response.

\subsection{Cost}

The work of an iteration is dominated by applying the Hamiltonian, local and nonlocal parts, to the $N_{\mathrm{b}}$ bands at each of the $n_k^{\mathrm{irr}}$ inequivalent $\vect{k}$ points, and by keeping those bands orthonormal. Applying the local part of the Hamiltonian to one band takes two FFTs, of order $N_{\mathrm{PW}}\log N_{\mathrm{PW}}$ each. The nonlocal projectors of the pseudopotential take of order $N_{\mathrm{PW}}N_{\mathrm{proj}}$, with $N_{\mathrm{proj}}$ the number of projectors in the cell. Keeping the bands orthonormal, with the diagonalization in the subspace they span, takes of order $N_{\mathrm{b}}^2N_{\mathrm{PW}} + N_{\mathrm{b}}^3$~\cite{payne1992iterative}. The cost per iteration is therefore
\begin{equation}
  t \propto n_k^{\mathrm{irr}}\bigl(N_{\mathrm{b}}N_{\mathrm{PW}}\log N_{\mathrm{PW}} + N_{\mathrm{b}}N_{\mathrm{PW}}N_{\mathrm{proj}} + N_{\mathrm{b}}^2N_{\mathrm{PW}}\bigr) .
  \label{eq:E.1-cost}
\end{equation}
Each term grows with the size of the cell. The number of bands and of projectors is proportional to the number of atoms $N_{\mathrm{at}}$, and so is $N_{\mathrm{PW}}$ at fixed cutoff, by \cref{eq:E.1-npw}. The last two terms therefore grow as $N_{\mathrm{at}}^3$ and dominate in large cells, while $n_k^{\mathrm{irr}}$ falls as $1/N_{\mathrm{at}}$ for a fixed sampling density. Raising the cutoff increases $N_{\mathrm{PW}}$ as $E_{\mathrm{cut}}^{3/2}$, which is why the 3d electrons of germanium, at \SI{80}{Ry} against \SI{60}{Ry}, cost a factor $(4/3)^{3/2} \approx 1.5$ in plane waves on top of the larger number of bands. Spin--orbit coupling doubles both the length of each state, now a two-component spinor, and the number of bands, and multiplies the $N_{\mathrm{b}}^2N_{\mathrm{PW}}$ term by eight. For a two-atom cell the cost is small: the self-consistent calculation of silicon at \SI{60}{Ry} on a $12^3$ grid converges in eight iterations and about four seconds on four cores. Cost becomes a constraint only for supercells, hybrids, the many empty states of a quasiparticle calculation, and the dense grids of response and electron--phonon calculations.

\subsection{Bands and effective masses}

Once the density is converged, it is held fixed and the Kohn--Sham Hamiltonian is diagonalized at any set of $\vect{k}$ points, on lines through the zone for a band structure, or on a fine grid around an extremum. This non-self-consistent step is cheap, because no iteration is needed, and it adds no approximation: the Kohn--Sham Hamiltonian is fixed by the converged density, and once the density is known the Hamiltonian can be evaluated at any $\vect{k}$.

An effective mass is the inverse curvature of a band at an extremum, $(m^{*-1})_{\alpha\beta} = \partial^2\varepsilon_{n\vect{k}}/\partial k_\alpha\partial k_\beta$ (\cref{sec:3.3}). It is computed by finite differences, by $\vect{k}\cdot\vect{p}$ theory or by Wannier interpolation. The three-point finite-difference formula along a direction,
\begin{equation}
  \frac{1}{m^*} \simeq \frac{\varepsilon(k_0 + h) - 2\varepsilon(k_0) + \varepsilon(k_0 - h)}{h^2} ,
  \label{eq:E.1-fd}
\end{equation}
has a truncation error $h^2\varepsilon''''/12$ from the next term of the Taylor series, and an error of about $4\delta\varepsilon/h^2$ if the eigenvalues carry an error $\delta\varepsilon$. A small step amplifies the noise of the eigenvalues, a large step samples the departure from a parabola, and the step must be chosen between the two. A least-squares parabola through several points around the extremum is more robust than \cref{eq:E.1-fd}, and varying the fitting range tests the truncation error and the noise error at once. The $\vect{k}\cdot\vect{p}$ route (\cref{sec:3.3}) uses perturbation theory with the computed velocity matrix elements, which gives the full tensor in one calculation at the extremum, provided enough bands are included and the velocities are corrected for the nonlocal pseudopotential (\cref{app:A.5}). Wannier interpolation (\cref{chap:27}) yields the band at any $\vect{k}$ and its derivatives analytically.

Computed masses inherit the errors of the functional. In $\vect{k}\cdot\vect{p}$ theory the curvature at an extremum comes from the coupling to the other bands at the same $\vect{k}$, each term weighted by the inverse energy separation. A mass set by a small direct gap, such as the conduction-band mass of GaAs at $\Gamma$, therefore follows the gap. A semilocal functional that puts the gap too low gives a mass that is too light, and a method that overestimates the gap gives one that is too heavy. Masses are checked against cyclotron resonance, and gaps against optical absorption or photoemission, before they are used in a rate.

\begin{workedexample}{Silicon from input to bands}\label{we:E-si-bands}
The representative calculation uses the PBE functional, the SG15 norm-conserving pseudopotential with four valence electrons, a cutoff of \SI{60}{Ry} (\SI{816}{eV}) for the states and \SI{240}{Ry} for the density, fixed occupations of the four valence bands, a $\Gamma$-centered $12^3$ grid (72 points by symmetry), scalar-relativistic potentials without spin--orbit coupling, and a stopping criterion of $10^{-12}$\,Ry on the change of the total energy. A relaxation of the lattice at this cutoff gives $a = \SI{5.477}{\angstrom}$, \SI{0.85}{\percent} above the measured room-temperature \SI{5.431}{\angstrom} (\cref{sec:3.1}; the zero-temperature value of \cref{tab:E-functionals}, with the zero-point expansion removed, is \SI{5.416}{\angstrom}), and \SI{0.2}{\percent} above the all-electron PBE value of that table, \SI{5.466}{\angstrom}. At the relaxed lattice constant the valence-band maximum is the threefold $\Gamma$ state, and the conduction-band minimum lies on the line from $\Gamma$ to X. A scan of that line in steps of $0.005\times2\pi/a$ puts the minimum at $k_c = (0.840 \pm 0.003)\times2\pi/a$, against the measured position near $0.85$ (\cref{sec:3.3}), and the indirect gap at $\SI{0.602}{eV}$, about half of the measured \SI{1.17}{eV}. The direct gap at $\Gamma$ is \SI{2.55}{eV}. At the measured lattice constant, with the same pseudopotential at \SI{40}{Ry}, where band energies are converged to about \SI{2}{meV}, the indirect gap is \SI{0.557}{eV}: the \SI{0.85}{\percent} expansion moves it by about \SI{45}{meV}.

The longitudinal mass of the conduction valley follows from the same scan. A parabola fitted to the nine eigenvalues within $\pm0.02\times2\pi/a$ of the minimum gives $m^*_\parallel = 0.953\,m_e$; widening the window to 13 and 17 points gives $0.954$ and $0.955$. The spread of the three masses, below one percent, is the combined truncation and rounding error of the fit, and $m^*_\parallel = 0.95\,m_e$ is the PBE value at this lattice constant. The measured value is $0.92\,m_e$~\cite{hensel1965cyclotron}. The PBE mass is \SI{4}{\percent} too heavy, while the gap is \SI{49}{\percent} too small: the longitudinal mass of the silicon valleys is insensitive to the error of the fundamental gap, and the functional error that matters for this material is the gap error itself.
\end{workedexample}

\subsection{Pitfalls}

A semilocal functional can get the order of the bands wrong, and no later correction of the gap alone repairs the wrong order. In germanium, PBE without spin--orbit coupling at the measured lattice constant places the $\Gamma$ conduction state \SI{0.047}{eV} above the valence-band maximum and the L state at \SI{0.191}{eV}, while in experiment the L valley is lowest, at \SI{0.74}{eV}~\cite[Table~I]{lucero2012improved}, and $\Gamma$ lies above it. At the larger PBE lattice constant the gap closes and the crystal becomes a metal. A rigid upward shift of the conduction bands keeps the wrong order, so any germanium band input needs bands that pass the check of \cref{app:E.14} (valley order, the L and $\Gamma$ gaps, masses at L and $\Gamma$), with spin--orbit coupling. A hybrid at the standard exchange fraction inverts $\Gamma$ and L even with spin--orbit coupling, and a shift of the 4s-like orbitals fitted to the L and $\Gamma$ gaps gets the masses wrong (\cref{app:E.14}). A quasiparticle calculation, tested by the same check, is the next option. The spin--orbit splitting of the valence-band maximum is \SI{0.296}{eV} in germanium against \SI{0.044}{eV} in silicon (\cref{sec:3.3}). The scalar-relativistic potentials of the worked example omit the splitting altogether; a calculation that includes it needs fully relativistic pseudopotentials and two-component spinors.

The gap error propagates into the response quantities: a gap that is too small gives too strong a screening. In germanium the overscreening becomes a failure: the near-metallic PBE state gives a dielectric constant of order $10^4$ that does not converge with the grid. When the computed dielectric constant fails to converge, either a corrected gap or the measured screening is used.

Smearing hides a metal. A semiconductor computed with smeared occupations, as some codes do by default, converges smoothly even when the functional has closed the gap, and the result looks normal unless the band structure is inspected. With fixed occupations, the same calculation struggles to converge, and the poor convergence is the useful symptom.
